
\subsubsection{Neural Network Calibration}

Guo et al. (2017) introduced temperature scaling for post-hoc calibration of image classifiers. The method learns a single temperature parameter T that scales logits before softmax computation. On CIFAR-100 and ImageNet, temperature scaling achieved 5-15\% ECE reduction. The method is theoretically grounded: it minimizes cross-entropy loss on held-out calibration data and is equivalent to maximum likelihood estimation under well-specified softmax parametrization.

Minderer et al. (2021) extended calibration analysis to vision transformers, demonstrating that modern architectures exhibit similar miscalibration patterns to convolutional networks. They confirmed that temperature scaling remains effective despite architectural changes.

Kadavath et al. (2022) observed qualitatively that large language models overestimate correctness probabilities on reasoning tasks. However, they did not quantify ECE or apply post-hoc calibration methods.

\subsubsection{Code Generation Evaluation}

Chen et al. (2021) introduced HumanEval with pass@k as the primary evaluation metric for code generation. Austin et al. (2021) released MBPP (Mostly Basic Python Problems) with 974 function-level tasks. Standard evaluation focuses exclusively on functional correctness (whether code passes tests) and does not assess confidence calibration.

\subsubsection{Positioning}

This work does not contribute new calibration methods or empirical findings about code generation models. Instead, it validates the experimental protocol for applying existing calibration methods \citep{guo2017calibration} to code generation tasks. The simulation establishes that the technical infrastructure is operational before committing to multi-hour production experiments.
